\item Dos bocinas idénticas separadas $10.0\text{ m}$ son conducidas por el mismo oscilador con una frecuencia $f = 21.5\text{ Hz}$ (figura P17.6) en una zona donde la rapidez del sonido es de $344\text{ m/s}$.
\begin{enumerate}
    \item Demuestre que un receptor en el punto A registra un mínimo de intensidad del sonido de las dos bocinas.
    \item Si el receptor se mueve en el plano de los altavoces, muestre la trayectoria que debe seguir para mantenerse en un mínimo la intensidad a lo largo de la hipérbola $9x^2 - 16y^2 = 144$.
    \item ¿Can el receptor permanecer en un mínimo y alejarse de las dos fuentes? Si es así, determine la forma limitante de la trayectoria que debe tomar. Si no es así, explique qué tanto se puede alejar.
\end{enumerate}